\clearpage\fancyhead[L]{\small VESTA.}\hypertarget{native-382-lower}{}\label{native-382-lower}
\begin{center}\textbf{VESTA.}\par\vspace{4pt}\hrule\vspace{4pt}\fontsize{10}{12}\selectfont Göttingische gelehrte Anzeigen. No. 129. pp. 1281\;..\;1284. 13 August 1814.\par\vspace{4pt}\hrule\vspace{4pt}\end{center}

\par\indent We have yet to give a brief report on the last opposition of \emph{Vesta}, in February of this year, and on some calculations prompted by it. At the observatory here, Professor \textsc{Gauss} made the following observations with the mural quadrant:

\begin{center}\fontsize{9}{11}\selectfont\setlength{\tabcolsep}{3pt}\renewcommand{\arraystretch}{1.15}\begin{tabular}{l|c|rrr|rrr}
1814 & \shortstack{Mean time\\in Göttingen} & \multicolumn{3}{c|}{\shortstack{Apparent right\\ascension}} & \multicolumn{3}{c}{\shortstack{Apparent\\declination}}\\\hline
Feb. 12 & $12^{\mathrm u}\;30^{\mathrm m}\;15^{\mathrm s}$ & $149^\circ$ & $56'$ & $48''3$ & $20^\circ$ & $47'$ & $3''3$ N.\\
14 & $12\;20\;26$ & 149 & 27 & 22.8 & 21 & 3 & 3.3\\
15 & $12\;15\;31$ & 149 & 12 & 35.5 & 21 & 10 & 55.7\\
20 & $11\;50\;55$ & 147 & 58 & 14.3 & 21 & 47 & 35.5\\
\end{tabular}\end{center}

{\rightskip=0pt plus 1em \par\indent Given the very slight illumination that the planet could bear in the telescope, which gathered little light, these observations can be credited with only moderate accuracy.\par}

\SourceLeaf{383}\Needspace{20\baselineskip}\par\indent The following observations were sent to Professor \textsc{Gauss} from other observatories:

\Needspace{14\baselineskip}\begin{center}\emph{Observations of Vesta at the Seeberg Observatory by Mr.} \textsc{Nicolai}.\end{center}
\begin{center}\fontsize{9}{11}\selectfont\setlength{\tabcolsep}{3pt}\renewcommand{\arraystretch}{1.15}\begin{tabular}{l|c|rrr|c}
1814 & \shortstack{Seeberg\\mean time} & \multicolumn{3}{c|}{\shortstack{Apparent right\\ascension}} & \shortstack{Apparent\\declination}\\\hline
Feb. 12 & $12^{\mathrm u}\;30^{\mathrm m}\;16^{\mathrm s}0$ & $149^\circ$ & $57'$ & $0''0$ & $\cdot\quad\cdot\quad\cdot$\\
15 & $12\;15\;31.6$ & 149 & 12 & 42.0 & $\cdot\quad\cdot\quad\cdot$\\
20 & $11\;50\;55.1$ & 147 & 58 & 15.9 & $21^\circ\;47'\;28''$ N.\\
25 & $11\;26\;25.1$ & 146 & 45 & 26.8 & $\cdot\quad\cdot\quad\cdot$\\
26 & $11\;21\;32.7$ & 146 & 31 & 17.7 & $\cdot\quad\cdot\quad\cdot$\\
28 & $11\;11\;50.7$ & 146 & 3 & 39.9 & $\cdot\quad\cdot\quad\cdot$\\
\end{tabular}\end{center}
\Needspace{17\baselineskip}\begin{center}\emph{Observations of Vesta at the Königsberg Observatory by Professor} \textsc{Bessel}.\end{center}
\begin{center}\fontsize{9}{11}\selectfont\setlength{\tabcolsep}{3pt}\renewcommand{\arraystretch}{1.15}\begin{tabular}{l|rrr|rrr|rrr}
1814 & \multicolumn{3}{c|}{\shortstack{Mean time\\in Königsberg}} & \multicolumn{3}{c|}{\shortstack{Apparent right\\ascension}} & \multicolumn{3}{c}{\shortstack{Apparent\\declination}} \\
\hline
Jan.\ 31  & $13^{\mathrm u}$ & $28^{\mathrm m}$ & $25^{\mathrm s}$ & $152^{\circ}$ & $40'$ & $16''8$ & $19^{\circ}$ & $6'$  & $38''4$ N. \\
Feb.\ 9  & 12 & 45 & 5  & 150 & 40 & 49.2 & 20 & 22 & 8.3 \\
13        & 12 & 25 & 30 & 149 & 42 & 44.3 & 20 & 54 & 38.2 \\
14        & 12 & 20 & 35 & 149 & 27 & 56.6 & 21 & 2  & 37.9 \\
19        & 11 & 55 & 58 & 148 & 13 & 33.3 & 21 & 40 & 19.8 \\
21        & 11 & 46 & 8  & 147 & 43 & 51.5 & 21 & 54 & 27.8 \\
22        & 11 & 41 & 13 & 147 & 29 & 9.1  & 22 & 1  & 9.7  \\
25        & 11 & 26 & 33 & 146 & 45 & 47.0 & 22 & 20 & 34.1 \\
26        & 11 & 21 & 40 & 146 & 31 & 36.8 & 22 & 26 & 34.1 \\
\end{tabular}\end{center}

\par\indent Professor \textsc{Bessel} determined the latitude of the new Königsberg Observatory to be $54^{\circ}\,42'\,50''31$ from 69 culminations of the Pole Star observed with the 2-foot \textsc{Cary} circle.

\Needspace{28\baselineskip}\par\indent Dr. \textsc{Gerling} in Cassel undertook to compare all these observations with the latest elements. To obtain better agreement, the epoch was revised. For 13 February at $12^{\mathrm u}$ mean time in Göttingen it was taken as $154^{\circ}\,54'\,8''39$. This calculation gave the following differences

\begin{center}\fontsize{10}{12}\selectfont\setlength{\tabcolsep}{6pt}\renewcommand{\arraystretch}{1.15}\begin{tabular}{l|c|c|l}
& \shortstack{in right\\ascension} & \shortstack{in\\declination} & Observer \\
\hline
Jan.\ 31 & $-11''15$ & $+4''44$  & Bessel \\
Feb.\ 9 & $-6.05$   & $+5.08$   & Bessel \\
12       & $-6.75$   & $\cdot\;\cdot\;\cdot\;\cdot$ & Nicolai \\
12       & $-16.45$  & $+10.42$  & Gauss \\
13       & $-3.16$   & $-0.11$   & Bessel \\
14       & $-3.13$   & $+5.62$   & Bessel \\
14       & $-11.23$  & $+16.76$  & Gauss \\
15       & $-8.95$   & $+22.86$  & Gauss \\
15       & $-0.46$   & $\cdot\;\cdot\;\cdot\;\cdot$ & Nicolai \\
19       & $+4.77$   & $+4.42$   & Bessel \\
20       & $+2.53$   & $+6.57$   & Nicolai \\
20       & $+2.98$   & $+0.21$   & Gauss \\
21       & $+2.61$   & $+6.97$   & Bessel \\
22       & $+2.48$   & $+0.93$   & Bessel \\
25       & $+3.60$   & $+5.05$   & Bessel \\
25       & $+6.28$   & $\cdot\;\cdot\;\cdot\;\cdot$ & Nicolai \\
26       & $+1.76$   & $+0.66$   & Bessel \\
26       & $+5.58$   & $\cdot\;\cdot\;\cdot\;\cdot$ & Nicolai \\
28       & $+7.92$   & $\cdot\;\cdot\;\cdot\;\cdot$ & Nicolai \\
\end{tabular}\end{center}

\par\indent Dr. \textsc{Gerling} then derived the opposition itself from these results and determined new elements that represent the four oppositions of 1810, 1811, 1812 and 1814.

\Needspace{12\baselineskip}\SourceLeaf{384}\emph{Fifth observed opposition of Vesta.}

\begin{center}13 February 1814. $9^{\mathrm u}\;5^{\mathrm m}\;40^{\mathrm s}$ Mean time in Göttingen\end{center}
\begin{center}\fontsize{10}{12}\selectfont\setlength{\tabcolsep}{2pt}\renewcommand{\arraystretch}{1.15}\begin{tabular}{lcr}
True longitude & $.\;.\;.\;.\;.\;.\;.\;.\;.\;.\;.\;.$ & $144^\circ\;34'\;54''47$\\
Geocentric latitude & $.\;.\;.\;.\;.\;.\;.\;.\;.\;.$ & $8\;2\;8.55\;\mathrm{N.}$\\
\end{tabular}\end{center}

\Needspace{16\baselineskip}\emph{New elements of Vesta.}

\begin{center}\fontsize{9}{11}\selectfont\setlength{\tabcolsep}{2pt}\renewcommand{\arraystretch}{1.15}\begin{tabular}{lcr}
Epoch of the mean longitude, 13 Feb. 1814. & & $12^{\mathrm u}$\\
in Göttingen & & $154^\circ\;55'\;27''83$\\
Longitude of perihelion & $.\;.\;.\;.\;.\;.\;.\;.\;.\;.\;.\;.\;.$ & $249\;38\;6.69$\\
Longitude of ascending node & $.\;.\;.\;.\;.\;.\;.\;.$ & $103\;11\;30.51$\\
\multicolumn{3}{c}{Both for the time of the epoch}\\
Inclination of orbit & $.\;.\;.\;.\;.\;.\;.\;.\;.\;.\;.\;.\;.$ & $7\;8\;16.01$\\
Eccentricity angle & $.\;.\;.\;.\;.\;.\;.\;.\;.\;.\;.\;.$ & $5\;8\;30.75$\\
Logarithm of mean distance & $.\;.\;.\;.\;.\;.\;.$ & $0.3731261$\\
Mean daily tropical motion & $.\;.\;.\;.\;.\;.$ & $977''95156$\\
\end{tabular}\end{center}

{\rightskip=0pt plus 1em \par\indent The same skilful astronomer is now computing an ephemeris for the next apparition of Vesta. The ephemeris will be published elsewhere shortly.\par}
